% Lösungen Einheit 4 von 4 – Nullstellen und Schnittpunkte berechnen (zusammenbau v0.1)
\einheitenkopf{Einheit 4 von 4 · Nullstellen und Schnittpunkte berechnen}

%% TODO Lösung der Erklärzeile 17g) fehlt in der Bank
\erg{17}{a) $x^2 - 2x = x + 4$ \quad b) $0 = (x - 1)^2 - 4$, $4 = (x - 1)^2$, $x - 1 = 2$ oder $x - 1 = -2$; $x_1 = 3$, $x_2 = -1$ \quad c) $0 = (x + 2)^2 - 9$, $9 = (x + 2)^2$, $x + 2 = 3$ oder $x + 2 = -3$; $x_1 = 1$, $x_2 = -5$ \quad d) $0 = (x - 4)^2 - 1$, $1 = (x - 4)^2$, $x - 4 = 1$ oder $x - 4 = -1$; $x_1 = 5$, $x_2 = 3$ \quad e) $0 = (x + 3)^2 - 25$, $25 = (x + 3)^2$, $x + 3 = 5$ oder $x + 3 = -5$; $x_1 = 2$, $x_2 = -8$ \quad f) $0 = (x - 5)^2 - 36$, $36 = (x - 5)^2$, $x - 5 = 6$ oder $x - 5 = -6$; $x_1 = 11$, $x_2 = -1$ \quad g) \ldots \quad h) Scheitel $S(2|3)$: keine Nullstelle – Scheitel über der $x$-Achse, nach oben geöffnet \quad i) $0 = x^2 - 2x - 8$; $x = 1 \pm \sqrt{1 + 8} = 1 \pm 3$; $x_1 = 4$, $x_2 = -2$ \quad j) $x = 2 \pm \sqrt{4 - 1} = 2 \pm \sqrt{3}$; $x_1 \approx 3{,}73$, $x_2 \approx 0{,}27$ \quad k) $0 = 3x^2 - 3x - 18$ $\mid : 3$: $0 = x^2 - x - 6$; $x = 0{,}5 \pm \sqrt{0{,}25 + 6} = 0{,}5 \pm 2{,}5$; $x_1 = 3$, $x_2 = -2$ \quad l) $x^2 + x + 1 = 7$, $x^2 + x - 6 = 0$; $x = -0{,}5 \pm \sqrt{0{,}25 + 6} = -0{,}5 \pm 2{,}5$; $x_1 = 2$, $x_2 = -3$ \quad m) $x^2 + x - 4 = 2x + 2$, $x^2 - x - 6 = 0$; $x = 0{,}5 \pm \sqrt{0{,}25 + 6} = 0{,}5 \pm 2{,}5$; $x_1 = 3$, $x_2 = -2$ \quad n) $(x - 2)^2 - 1 = x - 3$, $x^2 - 4x + 3 = x - 3$, $x^2 - 5x + 6 = 0$; $x = 2{,}5 \pm \sqrt{6{,}25 - 6} = 2{,}5 \pm 0{,}5$; $x_1 = 3$, $x_2 = 2$ \quad o) $x^2 - 4x + 1 = 2x - 4$, $x^2 - 6x + 5 = 0$; $x = 3 \pm \sqrt{9 - 5} = 3 \pm 2$; $x_1 = 5$, $x_2 = 1$; $g(5) = 6$, $g(1) = -2$; $S_1(5|6)$, $S_2(1|{-2})$ \quad p) $f(4) = 16 - 5 = 11$, $g(4) = 8 + 3 = 11$ – ja, $P$ liegt auf beiden Graphen. \quad q) z. B. $g(x) = -1$; jede waagerechte Gerade $g(x) = c$ mit $c$ kleiner als $1$, denn der Scheitel $S(2|1)$ ist der tiefste Punkt \quad r) $x^2 + 3x - 2 = x^2 - x + 6$, $4x = 8$, $x = 2$; $f(2) = 8$; $S(2|8)$ \quad s) $(x - 4)^2 - 6 = -2x + 5$, $x^2 - 8x + 10 = -2x + 5$, $x^2 - 6x + 5 = 0$; $x = 3 \pm \sqrt{9 - 5} = 3 \pm 2$; $x_1 = 5$, $x_2 = 1$; $g(5) = -5$, $g(1) = 3$; $S_1(5|{-5})$, $S_2(1|3)$ \quad t) $0 = x^2 - 10x + 22$; $x = 5 \pm \sqrt{25 - 22} = 5 \pm \sqrt{3}$; $x_1 \approx 6{,}73$, $x_2 \approx 3{,}27$; $N_1(6{,}73|0)$, $N_2(3{,}27|0)$ \quad u) $-x^2 + 4x + 2 = -10$, $-x^2 + 4x + 12 = 0$ $\mid \cdot (-1)$: $x^2 - 4x - 12 = 0$; $x = 2 \pm \sqrt{4 + 12} = 2 \pm 4$; $x_1 = 6$, $x_2 = -2$}
\erg{18}{a) Lina hat nur eine Lösung: auch $x + 1 = -5$ ist möglich; $x_1 = 4$, $x_2 = -6$ \quad b) Eine Nullstelle ist ein $x$ mit $f(x) = 0$; der Punkt $(x|0)$ hat den $y$-Wert null und liegt deshalb auf der $x$-Achse. \quad c) $S_1(1|{-2})$, $S_2(4|1)$; Probe: $f(4) = 1$ und $g(4) = 1$ \quad d) $0 = -0{,}1(x - 6)^2 + 4{,}9$, $(x - 6)^2 = 49$, $x = 13$ oder $x = -1$; sinnvoll ist $x = 13$: der Ball fliegt $13\,$m weit.}
